\documentclass[11pt]{article}
\usepackage[T1]{fontenc}
\usepackage{lmodern,microtype,amsmath,amssymb,amsthm,mathtools,booktabs,array}
\usepackage[margin=1in]{geometry}
\usepackage{flafter}
\usepackage{xurl}
\usepackage[colorlinks=true,allcolors=blue]{hyperref}
\hypersetup{pdftitle={Chain returns and lower bounds for superpermutations},pdfauthor={Jay Pantone}}
\newtheorem{theorem}{Theorem}[section]
\newtheorem{lemma}[theorem]{Lemma}
\newtheorem{proposition}[theorem]{Proposition}
\newtheorem{corollary}[theorem]{Corollary}
\DeclareMathOperator{\ord}{ord}
\DeclareMathOperator{\lcm}{lcm}
\DeclareMathOperator{\ex}{ex}
\newcommand{\CP}{\operatorname{CP}}
\newcommand{\ind}{\mathbf{1}}
\newcommand{\ceil}[1]{\left\lceil #1\right\rceil}
\newcommand{\floor}[1]{\left\lfloor #1\right\rfloor}
\title{Chain returns and lower bounds for superpermutations}
\author{Jay Pantone}
\date{October 2026}
\begin{document}
\maketitle

\begin{abstract}
A short superpermutation must make extensive use of large overlaps between
successive permutations. These overlaps organize the permutations into
runs, and earlier lower bounds control how efficiently such runs can be
joined. We strengthen this part of the argument. The patterns that attain
the earlier estimates cannot keep repeating: their action on the symbols
eventually brings them back to a rotation class already visited. In one
case, two repetitions already have an entire run in common. Counting the
necessary departures from these patterns gives an explicit lower bound
valid for every alphabet size at least eight. It improves the preceding
general formula at every size at least thirteen not divisible by three.
The leading correction above $K!+(K-1)!+(K-2)!$ remains asymptotic
to $2(K-4)!$. The exact formulas and their asymptotic comparisons have
formal proofs in Lean.
\end{abstract}

\noindent\textbf{AI disclosure.} The research presented here, the Lean
formalization, and this document were generated entirely by AI systems
under Jay Pantone's direction. The inherited mathematical work and
software retain their original credits. Independent AI review is not
human peer review.

\section{Introduction}

A \emph{superpermutation} on $K$ symbols is a word containing each of the
$K!$ permutations as a consecutive subword. For example, $123121321$
contains all six permutations on three symbols. Write $S(K)$ for the
length of a shortest superpermutation and put
\[
 F_3(K)=K!+(K-1)!+(K-2)!.
\]
The lower bounds studied here exceed $F_3(K)$ by a quantity asymptotic
to $2(K-4)!$. We improve the smaller terms by examining when the local
counting estimates can be tight. The resulting inequalities hold at every
alphabet size in their stated range.

The difficulty is that a pattern may be efficient for a few steps without
remaining efficient for very long. One can see the simplest instance
directly in a word. Appending the first symbol of a permutation produces
its next cyclic rotation, at a cost of only one letter. Repeating this
operation is inexpensive, but after $K$ steps it returns to the original
permutation. A covering word must eventually do something else.

The same phenomenon occurs further into the lower-bound argument.
After grouping rotations together, and then grouping those classes into
runs, the relevant objects are chains of almost full runs. The earlier
proof bounds the length of each chain and restricts which lengths may
follow one another. Its best average can nevertheless be attained by a
repeating sequence of chain lengths. Such a sequence does not determine
fresh permutations indefinitely. The corresponding relabeling of the
symbols has finite order, and when it returns, so does a previously
visited class.

This gives a way to improve the average. Each interruption of the
repeating pattern loses some of the length allowed by the old estimate.
If there is a uniform bound on the distance between interruptions, these
losses accumulate. There are two parts to making this observation useful:
finding a sufficiently short return, and proving that all other patterns
really do lose enough to pay for the interruptions. Most of the work in
this paper concerns the second part.

Our explicit bound is stated in Theorem~\ref{thm:main}. It applies for
every $K\ge8$, and incorporates the earlier refinement for $K$ divisible
by three. In each of the other two residue classes it strictly improves
the preceding general formula for every $K\ge13$. At $K=13$ and
$K=14$, for example, the increases are $12$ and $58$, respectively. The improvements hold uniformly
in the alphabet size; no separate search is used to prove an instance.

We use the geometric and forest lemmas of the lower-bound framework
retained in the accompanying proof library~\cite{pantone-framework}.
These lemmas reduce an arbitrary word to auxiliary paths of complete
rotation classes. The new estimates in this paper concern those paths.
The next section explains why counting deficits in such paths gives a
word-length bound. We then describe the chains that make the count
difficult. A fourteen-symbol example introduces the return argument
before we prove it in general. The explicit numerical formula is given
after these arguments, when its parameters have a combinatorial meaning.
The coordinate calculations and the endpoint count for the forest
reduction are given in the appendices.

\section{Overlaps, runs, and deficits}

To understand the quantities in the lower bound, we first describe how
the cheapest overlaps fit together. The \emph{weight} of a step from a
permutation $u$ to a different permutation $v$ is the number of letters
that must be appended to a word ending in $u$ to make it end in $v$.
Thus the weight is $K$ minus the length of the longest suffix of $u$
that is a prefix of $v$.
Ordering one chosen occurrence of each permutation in a covering word
produces a Hamiltonian path whose total weight is at most the word's
length minus $K$. Conversely, appending the required letters along such
a path produces a covering word. Hence $S(K)$ is $K$ plus the minimum
weight of a Hamiltonian path in this overlap graph.

A step has weight one precisely when it passes to the next cyclic
rotation. We call the $K$ rotations of a permutation a \emph{rotation
class}. Traversing all of them in cyclic order gives a \emph{complete
block}. The class of $1234$, for instance, consists of $1234,2341,3412,4123$.
Within this block, each of the three steps costs one. Repeating the next
step would revisit $1234$.

Our arguments take place in paths of complete blocks whose rotation
classes are pairwise distinct. The minimum-path and forest
reduction produces auxiliary paths of this kind from an arbitrary
superpermutation~\cite{pantone-framework}. Its exact form is recalled
in Appendix~\ref{app:forest}. A superpermutation as written
need not itself traverse whole classes at once. We shall return to this
reduction after seeing what can be counted in one of its paths.

It is convenient to put $n=K-1$ and use the alphabet $\{0,\ldots,n\}$.
Suppose a block is entered at $(x_1,\ldots,x_n,s)$. After traversing
the block, the exit is $(s,x_1,\ldots,x_n)$. A weight-two step to a different rotation class must
enter at $(x_2,\ldots,x_n,x_1,s)$.
The other ordering of the two remaining symbols lies in the class just
visited. Thus the candidate in a different class rotates the first $n$ symbols
and leaves the last symbol fixed. It may still be excluded if that class
has appeared earlier in the path.
Successive weight-two steps therefore arrange the blocks themselves
around a cycle of length $n$.

Figure~\ref{fig:run} shows this at $K=5$. Each row is one block, and
the four rows are joined by weight-two steps. Continuing would bring us
back to the first row, so a class-distinct path can use at most four
of these blocks consecutively.

\begin{figure}[htbp]
\centering
\renewcommand{\arraystretch}{1.25}
\begin{tabular}{ccccc}
$12340$ & $23401$ & $34012$ & $40123$ & $01234$\\
$23410$ & $34102$ & $41023$ & $10234$ & $02341$\\
$34120$ & $41203$ & $12034$ & $20341$ & $03412$\\
$41230$ & $12304$ & $23041$ & $30412$ & $04123$
\end{tabular}
\caption{A full run on five symbols. Read each row from left to right,
then pass to the next row. Steps within a row have weight one; steps
between rows have weight two. The same weight-two step from the last
row returns to the first.}
\label{fig:run}
\end{figure}

A maximal succession of blocks joined at weight two is a \emph{run}.
If it contains $L$ blocks, its \emph{deficit} is $n-L$. A full run has
deficit zero. For example, stopping Figure~\ref{fig:run} after the
third row gives a run of deficit one. A path with $r$ runs and total
deficit $D$ thus contains $nr-D$ blocks.

This identity explains the importance of deficits. They measure how much
of the capacity of the runs goes unused. If the total number of blocks
is fixed, a larger deficit forces more runs, and passing from one run
to the next costs an extra letter.

For a precise illustration, suppose a single path visits all $T=n!$
rotation classes and every boundary between runs has weight three.
There are $T$ blocks, $T-r$ weight-two boundaries, and $r-1$ weight-three
boundaries. Its word length is therefore
\begin{equation}\label{eq:counting-example}
 K+(K-1)T+2(T-r)+3(r-1)=K+(K+1)T-3+r.
\end{equation}
Ignoring deficits would give at least $T/n=(n-1)!$ runs. Substituting
this into \eqref{eq:counting-example} gives $F_3(K)+K-3$. Thus the
familiar factorial terms already appear in this simple count; the extra
contribution comes from proving that runs leave some capacity unused.

The quantity we need to bound is the number of runs. Since
$T=nr-D$, an estimate $D\ge\delta r-\beta$ gives
$T\le(n-\delta)r+\beta$. In other words, apart from the endpoint
allowance $\beta$, each run covers at most $n-\delta$ classes on average.
The denominator $n-\delta$ in the eventual lower bound has this meaning.

The forest reduction extends this accounting to arbitrary words. It
handles paths that break into components and boundaries of weight greater
than three, while keeping the same consequence of a deficit estimate.
The estimate must hold on every consecutive segment of runs, because
the proof applies it after cutting paths and adding certain entrance
blocks. This is the scope of the local arguments below.

\section{The chains that resist the count}

For the remainder of the paper, assume $n\ge7$. We now ask how a path
can have many runs while keeping its deficit small.
Restrict attention to consecutive runs joined at weight three; call such
a path a \emph{genuine light segment}. Its blocks still belong to distinct
rotation classes. Record the segment by its sequence of run deficits.
For example, an entry $2$ represents a run missing two blocks, while an
entry $0$ represents a full run.

The local overlap restrictions prevent us from inserting arbitrarily
many zeros between small positive deficits. Put $a=n-2$. A string of
full runs has length at most $a$, unless the whole segment consists
of full runs, in which case its length is at most $a+1$. The pattern
$(0,1,0)$ is impossible: a run missing one block cannot have full runs
on both sides. More generally, a factor $(0,1,d)$ has $d\ge2$.
Finally, if a gap of $a$ full runs separates two positive deficits,
those adjacent deficits sum to at least $a+2$. These are the earlier
local run restrictions~\cite{pantone-framework}.

Thus the first small positive deficit that can sit alone between full
runs is $2$. Two such entries can be separated by at most $a-1$ zeros.
The pattern that is most troublesome for the deficit count consists of
many singleton twos separated by these long gaps. Each two contributes
the same deficit, while the zeros increase the number of runs being
covered by the estimate.

This leads to the chains used in the proof. In a factor starting and
ending in zero, a positive cluster is \emph{regular} if it is the
singleton $2$. Join consecutive regular clusters across gaps of $a-1$
zeros to form a maximal \emph{chain}. Its length is its number of
regular clusters. A gap of $a-2$ zeros joins consecutive chains into a
\emph{path of chains}; shorter gaps and other positive clusters cut the
path. At $a=5$, for example, the factor
\[
 0,\ \underbrace{2,0,0,0,0,2}_{\text{chain of length 2}},\
 \underbrace{0,0,0}_{\text{shortened gap}},\
 \underbrace{2}_{\text{chain of length 1}},\ 0
\]
has two chains. These definitions retain the original adjacencies:
cutting out a nonregular cluster does not join its two neighboring gaps.

There are two ways to improve the count. One can limit the length of
a single chain, or show that several long chains cannot follow one
another too efficiently. The earlier proof already does both. We shall
strengthen the second estimate. The form we need is
\begin{equation}\label{eq:CP}
 \ell_1+\cdots+\ell_q\le p+\frac t2(q-1)
\end{equation}
for every path of $q$ chains. Here $p$ is an endpoint allowance, while
$t/2$ is the limiting average chain length permitted by the estimate.
Write $\CP(p,t)$ if \eqref{eq:CP} holds in every zero-ended factor of
every genuine light segment.

Each boundary between chains has one fewer full run than a link within
a chain. For a fixed number of regular clusters, more such boundaries
therefore allow fewer zero entries without changing the total deficit.
This is why shortening the permitted average chain length strengthens
the deficit estimate. The earlier scalar argument makes this precise: all the nonregular clusters
and shorter gaps already supply enough deficit to cover their part of
the count. We recall that argument in Appendix~\ref{app:scalar}.

\subsection{Following a chain also relabels the symbols}

The deficit sequence alone does not record which classes have been
visited. To recover some of this information, follow a chain in actual
permutations. Denote the two cells by
\[
 C=(2,0^{a-1}),\qquad B=(2,0^{a-2}),
\]
where an exponent means repetition. A chain of length $\ell$ followed
by a shortened gap has pattern $T_\ell=C^{\ell-1}B$. The final $B$
consists of the chain's last deficit-two run and its outgoing shortened
gap; this description is used only when that gap is present.

The cells are particularly useful because their internal steps are
forced. To see what this means, relabel a current exit as
$(0,1,\ldots,n)$. A weight-three step leads to $(3,\ldots,n,i,j,k)$,
where $(i,j,k)$ is one of the six orders of $(0,1,2)$. The local cell
calculation forces the order $210$. After a completed $C$, the same
choice is forced if the next run has at least two blocks; after $B$,
it is forced if the next run has at least three. The alternatives
collide with classes in the cell just traversed~\cite{pantone-framework}.

Consequently, once the first entry is fixed, each cell carries us to a
specified final exit. Relabeling that exit back to $(0,1,\ldots,n)$
puts the next cell into the same coordinates. Call this relabeling
the cell's \emph{normalization}. It acts on the visited classes as well
as on the exit. This last point is essential: resetting the coordinates
does not erase the earlier part of the path.

In one-line notation on $0,\ldots,n$, the two normalizations are
\begin{equation}\label{eq:maps}
 g=(0,5,4,6,\ldots,n,2,3,1),\qquad
 \rho=(0,1,3,4,\ldots,n,2),\qquad
 N(C)=g,\quad N(B)=\rho g.
\end{equation}
Products act rightmost first, so the normalization of $T_\ell$ is
$h_\ell=\rho g^\ell$. The maps compose in reverse chronological
order. If a succession of completed chains has normalization equal to
the identity, its actual following short-run start repeats the first
start class. That is impossible in a genuine segment.

Let $p_0$ be the order of $g$. The earlier return calculation gives
\[
\begin{array}{c|c}
n&p_0\\\hline
3m&2m\\
3m+1&m(2m+1)\\
3m+2&(m+1)(2m+1).
\end{array}
\]
For example, at $n=7$ the nonzero cycles of $g$ are $(1,5,2,4,7)$
and $(3,6)$, giving order $10$. A single chain has length at most
$p_0$. If there is more than one chain in the path, the earlier
endpoint collisions reduce every chain's maximum to $b=p_0-1$
\cite{pantone-framework}.

Knowing these maxima does not yet give the best average. A very long
chain can force its successor to be short. The next two sections examine
the patterns that attain the resulting averages, beginning with a case
where there is just one efficient pair.

\subsection{The inherited restrictions on consecutive chains}

We collect the local geometric facts used in the counts below. The
parameters $p_0$ in this statement are the orders just displayed. These
are restrictions on actual occupied classes, not on arbitrary integer
sequences with the same deficit sums.

\begin{lemma}[Local chain restrictions~\cite{pantone-framework}]
\label{lem:inherited-chains}
Every chain has length at most $p_0$. In a path of at least two chains,
every length is at most $b=p_0-1$. The following additional restrictions
hold whenever the indicated consecutive chains occur.

For $n=3m+1$ and $m\ge3$, in two consecutive chains, if the first is longer than $m$, then the
second has length at most $m$. That successor has length at most $m-1$ if it is
not the last chain of the path.

For $n=3m+2$ and $m\ge2$, put $v=2m(m+1)$. In two consecutive chains, if the first is longer than $v$, then the
second has length at most $v-1$. If the first chain has length
at least $v+2$ and its successor is nonterminal, the successor has
length at most $v-2$.

For $n=3m$ and $m\ge3$, four consecutive chains cannot have their
middle two lengths equal to $b$ and both outer lengths at least two.
For $n=7$, the chain inequality $\CP(10,10)$ holds.
\end{lemma}

The cell forcing rules and these chain restrictions are the inherited
geometric inputs to the present argument. Their proofs identify a class
in common between the relevant cell paths. Further occupied classes can
exclude more paths, but cannot undo such a collision. This is why the
same restrictions apply inside every factor to which the forest count
will be applied.

\section{Efficient pairs cannot keep repeating}\label{sec:pairs}

We begin with the fourteen-symbol case, because the improvement can be
seen before introducing its general parameters. Here $n=13$. A chain
has length at most $36$, and in a path with two or more chains its length
is at most $35$. The earlier local collision restrictions say more: a
chain longer than four forces its successor to have length at most four,
or at most three if that successor is not the last chain in the path.

Pair each nonterminal chain longer than four with its successor. The
pairs are disjoint, because the successor is too short to begin another
pair. All remaining chains, except possibly the last, have length at
most four. An interior pair has total length at most $38$, attained
only by $(35,3)$. Thus the earlier count allows an average of $19$
regular clusters per chain. The length count by itself would permit
$(35,3),(35,3),\ldots$ indefinitely.

Now follow what such a pair does to the symbols. Its normalization has
nonzero cycles
\[
 (1\ 11\ 7\ 3\ 9),\qquad(2\ 10\ 5\ 12\ 6),\qquad(4\ 13\ 8).
\]
After fifteen pairs the relabeling is the identity. If another chain
then starts, it starts in the same rotation class as the first one.
Consequently fifteen consecutive interior pairs cannot all be $(35,3)$.
At least one must have smaller total length, or the pairing must be
interrupted by a singleton. Either event loses at least one regular
cluster from the old allowance.

There are at most two chains per pair or singleton. Thus one lost unit
per fifteen such pieces predicts a reduction of $1/30$ in the mean
chain length. Equivalently, the parameter $t=38$ should decrease by
$1/15$. This is the source of the new value $569/15$. The argument
below proves this prediction while allowing for the final chain, whose
successor need not exist.

\subsection{Counting the interruptions}

Call a singleton or a chosen pair a \emph{tile}. If the old mean chain
length is $\mu$, a tile containing $k$ chains is allowed total length
$\mu k$. Its \emph{loss} is the difference between this allowance and
its actual total length. This loss counts regular clusters; the run
deficits introduced earlier count missing blocks.

The count is simple once the efficient tiles have been identified.
Mark every tile with positive loss. These marks divide the zero-loss
tiles into consecutive strings. A bound on the lengths of those strings
then forces enough marks. The only endpoint issue is that the last tile
may attain a larger length, and a return after it supplies no contradiction
unless a further chain actually starts. We therefore set that tile aside.

\begin{lemma}\label{lem:tiles}
Suppose a path of $q\ge2$ chains is partitioned into tiles of one or
two chains, and a tile with $k$ chains and total length $L$ has loss
$\mu k-L$. Suppose that before the last tile every loss is nonnegative,
every positive loss is at least $\varepsilon>0$, and every consecutive
string of zero-loss tiles has length at most $R-1$, where $R\ge1$.
If the last tile has loss at least $-c$, where $c\ge0$, then
\[
 \sum_{i=1}^q\ell_i\le
 \left(\mu-\frac{\varepsilon}{2R}\right)q+c+\varepsilon.
\]
\end{lemma}

\begin{proof}
Let $B_+$ be the number of marked tiles and $P$ their total loss,
excluding the last tile. Then $B_+\le P/\varepsilon$. There are at
most $B_++1$ intervening strings, each with at most $R-1$ zero-loss
tiles. Restoring the last tile, the total number of tiles is at most
$R(B_++1)$. Since each tile has at most two chains,
\[
 q\le2R(P/\varepsilon+1),\qquad
 P\ge\frac{\varepsilon q}{2R}-\varepsilon.
\]
The total chain length is at most $\mu q-P+c$, giving the result.
\end{proof}

\subsection{The case \texorpdfstring{$n=3m+1$}{n=3m+1}}

The fourteen-symbol example is typical of this residue class. Assume
$m\ge3$, put $p_0=m(2m+1)$ and $b=p_0-1$, and pair every nonterminal
chain longer than $m$ with its successor. Lemma~\ref{lem:inherited-chains} bounds that successor by $m$, and
by $m-1$ when it is nonterminal. Hence an interior pair has total
length at most $t_0=p_0+m-2$, with equality only for $(b,m-1)$.

Take $\mu=t_0/2$. An interior singleton has length at most $m$, so
its loss is at least $\mu-m\ge1$. All positive interior losses are
integers, since $t_0=2m^2+2m-2$ is even. Thus a tile either is the
efficient pair $(b,m-1)$ or loses at least one unit.

The normalization of that pair is
\begin{equation}\label{eq:efficient-map}
 M_1=h_{m-1}h_b=\rho g^{m-1}\rho g^{-1}.
\end{equation}
Let $R_1$ be its order. As in the example, $R_1$ consecutive efficient
interior tiles would return to the first start class. The next start
is forced because a deficit-two run has $n-2\ge8$ blocks, above the
three-block threshold after $B$.

Merely knowing that $M_1$ has finite order would give a weak quantitative
estimate. Here its cycles can be calculated explicitly. Their lengths
are listed in Table~\ref{tab:cycles}, and the calculation in
Appendix~\ref{app:cycles} proves $R_1\le n^2/4$. The example $n=13$
has cycle lengths $5,5,3$, so $R_1=15$.

It remains to account for the last tile. A final singleton can have
length $b$, giving loss $-(p_0-m)/2$. A final pair can have total
length $b+m=t_0+1$, giving loss $-1$. Thus the single allowance
$c=(p_0-m)/2$ covers both possibilities. Apply Lemma~\ref{lem:tiles}
with $\varepsilon=1$. Since $p_0-\mu=c+1$, it yields
\[
 \sum_i\ell_i\le
 \left(\mu-\frac1{2R_1}\right)q+c+1
 \le p_0+\frac12\left(t_0-\frac1{R_1}\right)(q-1).
\]
The case of a single chain already has the stronger bound $p_0$.
We have therefore proved
\begin{equation}\label{eq:CP1}
 \CP\left(p_0,t_0-\frac1{R_1}\right).
\end{equation}
At fourteen symbols, the final-tile allowance is $16$. It remains an
endpoint allowance; it does not cancel the losses that accumulate along
a long path.

\subsection{When the interiors collide first}

For $n=3m+2$, following the first class of each tile misses an earlier
collision. It is useful to recall the full run in Figure~\ref{fig:run}.
Delete the symbol zero from each of its four rotation classes. All four
give the same cyclic order $1234$; they differ only in the gap into which
zero was inserted. In general, a full run with distinguished symbol zero
visits all $n$ such insertions into one cyclic order of the other symbols.
This family of classes is called a \emph{two-loop}.

Consequently two full runs with the same nonzero cyclic order repeat
every class, even if the starts of the larger pieces containing them
are different. This is what happens in the efficient pattern for the
present residue. Put $p_0=(m+1)(2m+1)$, $v=2m(m+1)$, $b=p_0-1$,
and $c=v-2$.

\begin{lemma}\label{lem:two-pairs}
For $m\ge2$, the pattern $T_bT_cT_bT_c$ cannot occur in a genuine
light segment.
\end{lemma}

The repeated runs lie inside the two shorter chains. To locate them,
look at the $C$ cells of each $T_c=C^{c-1}B$. In the first copy choose
the cell at distance $m$ from the end of this $C$ portion; in the second
choose the cell at distance $m$ from its beginning, as in
Figure~\ref{fig:collision}. The last full run of the first chosen cell
and the first full run of the second chosen cell have the same nonzero
cyclic order. Appendix~\ref{app:full-collision} checks this equality by
following the relabelings. Both runs lie strictly before their chains'
final $B$ cells, so the collision occurs within the displayed pattern.

\begin{figure}[htbp]
\centering
\renewcommand{\arraystretch}{1.6}
\begin{tabular}{r l l}
first pair & $T_b\quad C^{c-2-m}\,\boxed{C}\,C^mB$
 & last full run in the box\\
second pair & $T_b\quad C^m\,\boxed{C}\,C^{c-2-m}B$
 & first full run in the box
\end{tabular}
\caption{The two full runs that coincide. The boxes are at equal depths
from opposite ends of the $C$ portions of the shorter chains. Their
cyclic orders are compared in the coordinates of the whole path.}
\label{fig:collision}
\end{figure}

We can now use the same loss count. Lemma~\ref{lem:inherited-chains} bounds the successor of a chain
longer than $v$ by $v-1$. If the first length is at least $v+2$ and
the successor is nonterminal, the stronger bound is $v-2$.
Pair every nonterminal chain longer than $v$ with its successor, and
put $t_0=b+v-2=4m^2+5m-2$ and $\mu=t_0/2$.

For $m>2$, the only interior pair attaining $t_0$ is $(b,v-2)$.
A pair beginning with $v+1$ loses at least $m-2$, and every other
positive pair loss is an integer. A singleton loses at least
$(m-2)/2$. Thus every positive interior loss is at least
$\varepsilon=\min\{1,(m-2)/2\}$. Lemma~\ref{lem:two-pairs}
excludes two successive efficient tiles, so take $R=2$.

A final singleton has loss at least $-(m+2)/2$, and a final pair has
loss at least $-1$. Lemma~\ref{lem:tiles}, together with
$p_0-\mu=(m+2)/2+1$ and $\varepsilon\le1$, gives
\begin{equation}\label{eq:CP2}
 \CP(p_0,t_0-\varepsilon/2),\qquad m>2.
\end{equation}
The decrement is $1/4$ at $m=3$ and $1/2$ thereafter.

At $m=2$ there are three efficient tiles: $(12)$, $(13,11)$ and
$(14,10)$. Their normalizations lie in a group of order $12$, so
twelve consecutive efficient interior tiles would have two equal
prefix normalizations and hence a repeated start class. Positive
losses are integers and the final allowance is two. The same lemma
gives $\CP(15,287/12)$. The small group calculation is also included
in Appendix~\ref{app:full-collision}.

All these collisions occur within the path under consideration. Thus
the counts apply in every zero-ended factor, as required by $\CP$;
no completion of such a factor to a superpermutation is assumed.

\section{When the efficient pattern alternates}

Assume $n=3m$ with $m\ge3$. This case has a different efficient pattern. Here the period is
$p_0=2m$, so every chain in a nontrivial path has length at most
$b=p_0-1$. Lemma~\ref{lem:inherited-chains} forbids four consecutive chains
whose middle two both have length $b$ and whose outer two have
length at least two.

Temporarily suppose that no chain has length one. Away from the ends of
the path, two chains of maximum length cannot then be adjacent. The
largest average allowed by this restriction comes from alternating $b$
and $b-1$. For example, at $n=12$ the candidate sequence is
$7,6,7,6,\ldots$, with mean $13/2$. The old parameter is accordingly
$t_0=2b-1=2p_0-3$.

As before, this efficient repetition returns to an occupied class. The
additional issue is how to count departures from it. A chain of length
one can allow a neighboring pair of maximum lengths that would otherwise
be forbidden. Such a short chain loses many units of length, and some of
that loss must pay for the exception it creates. Counting only the number
of nonalternating positions would discard this information.

\subsection{The return of an alternating pair}

The normalization of one pair is $h_{b-1}h_b$. It has an especially
simple description after changing coordinates: advance seven places
backwards around a circle, then interchange two neighboring symbols.
A transposition either splits one cycle into two or joins two cycles
into one. This gives the order without following a long word of cells.

Define
\begin{equation}\label{eq:R0}
 R_0=\begin{cases}
 2n/7,&7\mid n,\\
 u(n-u),&7\nmid n,
 \end{cases}
\end{equation}
where in the second case $1\le u<n$ and $7u\equiv1\pmod n$.
For $n=12$, the inverse is $u=7$, giving $R_0=35$.

\begin{lemma}\label{lem:alternating-return}
An alternating factor of $2R_0+1$ chains, with lengths $b,b-1$ in
either order, cannot occur in a genuine light segment. Moreover,
$R_0\le n^2/4$.
\end{lemma}

\begin{proof}
Let $s$ be the successor permutation of the circle $(1,3,\ldots,n,2)$,
and put $\tau=(1\ 3)$. The permutations in \eqref{eq:maps} give
$\rho=\tau s$ and $g=\tau s^3$. Since $g$ has order $p_0$,
\[
 M_0=h_{b-1}h_b=\rho g^{-2}\rho g^{-1}
     =\tau s^{-2}\tau s^{-5}\tau.
\]
Conjugation by $s^2\tau$ gives $\tau s^{-7}$, as described above.
If $7\nmid n$, the map $s^{-7}$ is an $n$-cycle. The transposition
cuts it into arcs of lengths $u$ and $n-u$. These lengths are coprime,
so the order is $u(n-u)\le n^2/4$.
If $7\mid n$, the map $s^{-7}$ has seven cycles of length $n/7$.
The transposition joins two and leaves the other five, giving order
$2n/7\le n^2/4$.

Reversing the pair gives a conjugate map with the same order. In an
alternating factor of $2R_0+1$ chains, the first $2R_0$ have outgoing
shortened gaps. Their $R_0$ pairs have combined normalization equal
to the identity, so the following short-run start repeats the first
start class. The post-$B$ forcing rule applies because each such run
has $n-2\ge7$ blocks.
\end{proof}

\subsection{Paying for the breaks in alternation}

We first describe the count. Measure each chain's loss from its maximum
length $b$. The two efficient lengths $b,b-1$ then become the binary
symbols $0,1$, and the efficient pattern becomes an alternating binary
string. We shall remove larger losses, cut at adjacent equal binary
symbols, and apply Lemma~\ref{lem:alternating-return} to the alternating
strings that remain.

In the interior of the path there are two kinds of equal adjacency.
A pair $11$ already contains
extra loss: both chains have fallen below the maximum. A pair $00$ has
no such loss, but the four-chain restriction requires a neighboring
chain of length one. That chain supplies the payment. The following
proof keeps track of these two sources separately.

\begin{proposition}\label{prop:CP0}
For $n=3m$, $m\ge3$, we have
\[
 \CP\left(p_0+\frac12+\frac1{4R_0},\,
 t_0-\frac1{2R_0}\right).
\]
\end{proposition}

\begin{proof}
Consider a path of $q\ge2$ chains. Write $z_i=b-\ell_i$ for the loss
of its $i$th chain, and let $w_i$ indicate whether that chain has
length one. For each four-chain window, at least one of the middle
chains is below maximum length or one of the outer chains has length
one. Thus, for $2\le i\le q-2$,
\[
 f_i:=\ind_{z_i>0}+\ind_{z_{i+1}>0}+w_{i-1}+w_{i+2}\ge1.
\]
Think of each summand as a witness to this local restriction.
An ordinary nonmaximum chain supplies one marker, and a chain of length
one supplies two: it can witness the restriction in a middle position
or in an outer position. Accordingly, put $v_i=\ind_{z_i>0}+w_i$.
Each marker occurs in at most two of the windows counted by $\sum f_i$.
Also $v_i\le z_i$: a chain of length one loses $b-1\ge4$ units,
more than enough to supply its two markers.

There are $q-3$ windows when $q\ge3$. This suggests measuring the
remaining loss by $E=2\sum_i z_i-(q-3)$, after subtracting their
mandatory witnesses. Let $A=\sum_{i=2}^{q-2}(f_i-1)$ count the extra
witnesses. Then
\[
 E=2\sum_i(z_i-v_i)
   +\left(2\sum_i v_i-\sum_{i=2}^{q-2}f_i\right)+A.
\]
All three terms are nonnegative. A nonmaximum chain other than a chain
of length one spends one loss unit on its marker. If its loss is at
least two, at least one unit remains. A chain of length one spends two
of its $b-1$ units, leaving $b-3$. Consequently, if $B_0$ counts the
indices with $z_i\ge2$ and $\ell_i\ne1$, and $W$ counts those with
$\ell_i=1$, then
\begin{equation}\label{eq:defect-payment}
 E\ge2B_0+2(b-3)W+A.
\end{equation}
This is the reserve available to pay for cuts.

Now restrict to the interior indices $2,\ldots,q-1$ and remove those
counted by $B_0$ or $W$, leaving breaks where they stood. The surviving
strings are binary. Cut at every $00$ or $11$ adjacency. Each $11$
contributes at least one extra witness to $A$. Each $00$ needs a
neighboring chain of length one, and a single such chain can explain
at most two $00$ adjacencies. Hence the total number of alternating
strings is at most $B_0+3W+A+1$: a removed chain creates one break,
and a chain of length one may account for two further cuts.

Every remaining string has length at most $2R_0$. Counting the removed
indices as well gives
\begin{equation}\label{eq:alternating-count}
 q-2\le B_0+W+2R_0(B_0+3W+A+1)\le2R_0(E+1).
\end{equation}
For the last inequality, \eqref{eq:defect-payment} pays for each type
of break separately; $b\ge5$ and $R_0\ge1$ suffice for the coefficients.
At $q=3$ the window sums are empty and the argument still applies.
At $q=2$, the identity $E=2\sum_i z_i+1\ge1$ gives
\eqref{eq:alternating-count} directly.

Finally let $\mu=t_0/2=b-1/2$. The total length is
\begin{equation}\label{eq:loss-length}
 \sum_i\ell_i=\mu q+\frac32-\frac E2.
\end{equation}
Inserting \eqref{eq:alternating-count} into \eqref{eq:loss-length}
gives
\[
 \sum_i\ell_i\le
 \left(\mu-\frac1{4R_0}\right)q+2+\frac1{2R_0}
 =p_0+\frac12+\frac1{4R_0}
  +\frac12\left(t_0-\frac1{2R_0}\right)(q-1).
\]
For a single chain, the earlier bound $\ell_1\le p_0$ suffices.
\end{proof}

The average has decreased, but the endpoint allowance has increased
slightly. The count explains both changes: alternating pieces have bounded
length, while the two ends can leave incomplete pieces. Only the first
cost accumulates with the number of chains.
The same proof permits any real upper bound $U\ge R_0$ in place of
$R_0$, because every alternating string has length at most $2R_0\le2U$.
In particular, the simpler choice $U=n^2/4$ is valid.

\section{The resulting lower bound}

We have obtained the chain estimates needed by the deficit count. We now
apply the earlier word reduction and give the parameters together. This
is the same capacity calculation illustrated in \eqref{eq:counting-example},
with the forest theorem accounting for the components and their endpoints.

For $p\ge0$ and $t>0$, define
\begin{equation}\label{eq:bound}
 \mathcal B_K(p,t)=F_3(K)+K-3+
 \ceil{\frac{2(K-2)!-2(K-2)-4p/t}
 {K^2-4K+1-2(K-1)/t}}.
\end{equation}
In each application below the fraction is positive. The preceding
sections have supplied $p$ and $t$ by bounding the total length of a path
of chains. The following proposition states precisely how that estimate
enters the word bound.

\begin{proposition}[Chain-to-word transfer]\label{prop:transfer}
Let $n=K-1\ge7$. Suppose $p,t$ are real numbers satisfying
\[
 0\le p\le t,\qquad t>0,\qquad t(n-4)\ge4p+6.
\]
If $\CP(p,t)$ holds, then $S(K)\ge\mathcal B_K(p,t)$.
\end{proposition}


The proof combines the scalar deficit estimate with the hereditary
forest theorem. For reference, it is included with the scalar argument
in Appendices~\ref{app:scalar} and~\ref{app:forest}. Its quantifier over all genuine segments
is what permits the estimate to be used on the paths produced by the
reduction from an arbitrary word.

The two return orders in the next theorem are $R_0$ from
\eqref{eq:R0} and $R_1=\ord(M_1)$ from \eqref{eq:efficient-map}.
Their exact values can make the bound stronger at particular sizes.
Replacing them by their common upper bound $n^2/4$ gives a simpler
formula consisting of rational functions on the three residue classes.

\begin{theorem}\label{thm:main}
For every $K\ge8$, we have $S(K)\ge\mathcal B_K(p,t)$ with the following
parameters. If $K\ge13$, write $n=K-1=3m+r$, where $0\le r\le2$, and use
\[
\renewcommand{\arraystretch}{1.5}
\begin{array}{c|c|c}
r&p&t\\\hline
0&2m+\dfrac12+\dfrac1{4R_0}&4m-3-\dfrac1{2R_0}\\
1&m(2m+1)&2m^2+2m-2-\dfrac1{R_1}\\
2&(m+1)(2m+1)&4m^2+5m-\dfrac52
\end{array}
\]
For $8\le K\le12$, use
\[
\renewcommand{\arraystretch}{1.2}
\begin{array}{c|ccccc}
K&8&9&10&11&12\\\hline
p&10&15&521/80&21&28\\
t&46/3&287/12&359/40&263/12&195/4.
\end{array}
\]
In the general table, either return order may instead be replaced by
$n^2/4$.
\end{theorem}


\begin{proof}[Proof of Theorem~\ref{thm:main}]
The chain inequalities are \eqref{eq:CP1}, \eqref{eq:CP2}, and
Proposition~\ref{prop:CP0}. At $K=8$, the earlier theorem gives
$\CP(10,10)$, which implies $\CP(10,46/3)$.
The numerical hypotheses of Proposition~\ref{prop:transfer} are verified
in Appendix~\ref{app:numerical}, including all five initial sizes.
That proposition gives the word bound. The counts remain valid when a
return order is replaced by any larger real number, proving the final
assertion as well.
\end{proof}

Write $B(K)$ for the exact-order formula in Theorem~\ref{thm:main}
and $Q(K)$ for the quadratic-order formula. For $8\le K\le12$ these
use the same initial rows. Every occurrence of a difference between
these formulas below is an ordinary difference of integers.

\subsection{The contribution of the return orders}

To isolate the return-order improvements, define $O(K)$ by the same
initial rows for $K<13$ and, for $K\ge13$, by
$O(K)=\mathcal B_K(p,t)$ with the parameters
\[
\begin{array}{c|c|c}
n=K-1 & p & t\\\hline
3m & 2m & 4m-3\\
3m+1 & m(2m+1) & 2m^2+2m-2\\
3m+2 & (m+1)(2m+1) & 4m^2+5m-5/2.
\end{array}
\]
Thus $O$ already includes the collision between two efficient pairs in
Section~\ref{sec:pairs}. The remaining differences come from the
orders $R_0$ and $R_1$.

\begin{proposition}\label{prop:gains}
For every $K\ge13$ with $K\equiv1,2\pmod3$, we have
$B(K)>O(K)$. Along these two residue
classes, respectively,
\[
 \liminf_{K\to\infty}
 \frac{B(K)-O(K)}{(K-9)!}
 \ge\frac92,
 \qquad
 \liminf_{K\to\infty}
 \frac{B(K)-O(K)}{(K-11)!}
 \ge324.
\]
Replacing $B$ by $Q$, both normalized differences converge to the
displayed constants.
\end{proposition}

\begin{proof}
Before taking ceilings, the exact-order improvements in the two cases are
asymptotic to
\[
 \frac{9n!}{8R_0n^6},\qquad \frac{81n!}{R_1n^8},
\]
respectively. Appendix~\ref{app:numerical} gives the exact subtractions
and verifies that the real improvement exceeds one from $n=12$ and
$n=13$, respectively. The bound $R_i\le n^2/4$ gives the two lower
limits, with equality in the limit for the quadratic choices.
Taking ceilings changes each difference by less than one.
\end{proof}

The exact orders sometimes give larger gains: for instance,
$R_0=2n/7$ when $7\mid n$, and $R_1=n$ in two residue classes of
$m$. This is why Proposition~\ref{prop:gains} states lower limits
for the strongest formula. At $K=13,14$, its gains are $12$ and
$58$; at $K=16,17$, they are $9\,528$ and $16\,920$.

\subsection{Comparison with the all-size reference bound}

Halfar~\cite{halfar} proves the following all-size formula, which he
attributes to Fritsch:
\begin{equation}\label{eq:H}
 H(K)=F_3(K)+K-3+
 \ceil{\frac{2((K-2)!-(K-2))}{K^2-4K+1}},\qquad K\ge7.
\end{equation}
This is Theorem B of the proof package announced on 6 October
2026~\cite{halfar-announcement}. His Theorem C gives stronger bounds
from searched parameters at individual sizes. The comparison below
is with the displayed formula $H$, not those finite bounds.

The earlier chain theorem already improves the all-size formula
\eqref{eq:H} at a larger order than the refinements above. We record
the comparison to distinguish these two contributions.

\begin{corollary}\label{cor:comparison}
For either $A=B$ or $A=Q$, we have
\[
 A(K)=F_3(K)+(2+o(1))(K-4)!,
\]
and, as $K\to\infty$ through the indicated residue class,
\[
 A(K)-H(K)=
 \begin{cases}
 (3+o(1))(K-6)!,&K\equiv1\pmod3,\\
 (18+o(1))(K-7)!,&K\equiv2\pmod3,\\
 (9+o(1))(K-7)!,&K\equiv0\pmod3.
 \end{cases}
\]
\end{corollary}

\begin{proof}
For both parameter choices, $p/t$ is bounded and $t\ge1$ for all
sufficiently large $K$. The numerator inside the ceiling in
\eqref{eq:bound} is therefore $2(K-2)!+O(K)$ and its denominator is
$K^2(1+O(1/K))$. Since $(K-2)!/K^2\sim(K-4)!$, the claimed leading
term follows. The additive term $K-3$ and the rounding error are
$o((K-4)!)$.
With $\kappa=n^2-2n-2$, the difference of the real expressions
inside the ceilings in \eqref{eq:bound} and \eqref{eq:H} is exactly
\[
 \frac{4[n!-n(n-1)-p\kappa]}{t\kappa(\kappa-2n/t)}.
\]
In the three cases $n\equiv0,1,2\pmod3$, the parameter $t$ is
asymptotic to $4n/3$, $2n^2/9$, and $4n^2/9$, respectively.
Substitution gives the three constants. Rounding changes each
difference by less than one.
\end{proof}

\begin{corollary}\label{cor:word-asymptotic}
For every real $\varepsilon>0$, there is $K_0$ such that every
superpermutation on $K\ge K_0$ symbols has length at least
\[
 F_3(K)+(2-\varepsilon)(K-4)!.
\]
\end{corollary}

\begin{proof}
The normalized correction of $B(K)$ tends to $2$, and every covering
word has length at least $B(K)$ by Theorem~\ref{thm:main}.
\end{proof}

The conclusion concerns a lower bound, not an asymptotic formula for
$S(K)$. In particular, the present argument does not produce a positive
coefficient of $(K-3)!$. It strengthens the count of the cheapest local
patterns while retaining the same segment-to-word reduction.

\appendix
\section{The return order when \texorpdfstring{$n=3m+1$}{n=3m+1}}\label{app:cycles}

We calculate the order of the map $M_1$ in \eqref{eq:efficient-map} for $m\ge3$.
\begin{table}[htbp]
\centering
\renewcommand{\arraystretch}{1.3}
\begin{tabular}{c l}
\toprule
$m\bmod6$ & Cycle lengths defining $R_1$\\
\midrule
$0,2$ & $3m+1$\\
$1$ & $(2m+1)/3,\ (7m+2)/3$\\
$3$ & $(7m+3)/6,\ (11m+3)/6$\\
$4$ & $(2m+1)/3,\ (7m+2)/6,\ (7m+2)/6$\\
$5$ & $(7m+1)/6,\ (11m+5)/6$\\
\bottomrule
\end{tabular}
\caption{The nonzero cycles of the return map when $n=3m+1$.
The order $R_1$ is their least common multiple.}
\label{tab:cycles}
\end{table}

The nonzero cycles of $g$ are
\[
 (3,6,\ldots,3m),\qquad
 (1,5,8,\ldots,3m-1,2,4,7,\ldots,3m+1).
\]
The first has length $m$ and the second length $2m+1$. To evaluate
$g^{m-1}$ on the second, index its letters by
\[
 B_0=1,\quad B_j=3j+2\ (1\le j<m),\quad B_m=2,\quad
 B_{m+j}=3j+1\ (1\le j\le m).
\]
The map $g$ adds one to these indices modulo $2m+1$. Also
\[
 (\rho g^{-1})(1)=2,\quad
 (\rho g^{-1})(2)=3m,\quad
 (\rho g^{-1})(3)=3m+1,\quad
 (\rho g^{-1})(4)=3,\quad
 (\rho g^{-1})(5)=1,
\]
and $(\rho g^{-1})(x)=x-2$ for $6\le x\le3m+1$.
Applying $g^{m-1}$ and then $\rho$ gives
\[
\begin{array}{c|rrrrrrrrr}
x&1&2&3&4&5&6&7&8&9\\\hline
M_1(x)&3m-1&3m-2&3m-3&3m+1&3m&2&3&4&1,
\end{array}
\]
and, for $10\le x\le3m+1$,
\begin{equation}\label{eq:descent}
 M_1(x)=\begin{cases}
 x-6,&x\equiv0\pmod3,\\
 x-5,&x\equiv1\pmod3,\\
 x-4,&x\equiv2\pmod3.
 \end{cases}
\end{equation}
For example, when $x=3i\ge6$, its image under $\rho g^{-1}$ is
$B_{m+i-1}$. Adding $m-1$ to this index gives $2m$ at $i=2$, zero
at $i=3$, and $i-3$ for $i\ge4$. These are the exceptional values
$M_1(6)=2$, $M_1(9)=1$, and the first line of \eqref{eq:descent}.
The other two residues follow by the same index substitution.

Every orbit outside $\{1,2,3,4,5\}$ decreases until it reaches that
set. Thus every cycle contains one of these five \emph{anchors}.
The descent times are as follows; the first jump out of an anchor is
not included.
\[
\renewcommand{\arraystretch}{1.15}
\begin{array}{c|c|c}
\text{Starting value}&\text{Anchor reached}&\text{Steps}\\\hline
6j\ (j\ge1)&2&j\\
6j+3\ (j\ge1)&1&j\\
9j+1\ (j\ge1)&5&2j-1\\
9j+4\ (j\ge0)&4&2j\\
9j+7\ (j\ge0)&3&2j+1\\
9j+2\ (j\ge1)&3&2j\\
9j+5\ (j\ge0)&5&2j\\
9j+8\ (j\ge0)&4&2j+1
\end{array}
\]
For multiples of three, this is repeated subtraction of six, ending
at $6\mapsto2$ or $9\mapsto1$. In the other residues, two steps
subtract nine. The low values $4,5,7,8,10$ give the stated endpoints.

Write $m=6h+r$. Applying the descent times to the five initial jumps
gives the following first-return data, in anchor order $1,2,3,4,5$.
\[
\renewcommand{\arraystretch}{1.2}
\begin{array}{c|c|c}
r&\text{Destinations}&\text{Arc lengths}\\\hline
0&(4,3,1,5,2)&(4h,4h,3h,4h,3h+1)\\
1&(3,5,2,4,1)&(4h+1,4h,3h+1,4h+1,3h+1)\\
2&(5,4,1,3,2)&(4h+1,4h+1,3h+1,4h+2,3h+2)\\
3&(4,3,2,5,1)&(4h+2,4h+2,3h+2,4h+2,3h+2)\\
4&(3,5,1,4,2)&(4h+3,4h+2,3h+2,4h+3,3h+3)\\
5&(5,4,2,3,1)&(4h+3,4h+3,3h+3,4h+4,3h+3).
\end{array}
\]
Summing the arc lengths around each anchor cycle gives exactly
Table~\ref{tab:cycles}. There are no further cycles, since every
nonanchor reaches an anchor.

For a single cycle, the order is $n\le n^2/4$. For two cycles of
lengths $u,v$, it is at most $uv\le n^2/4$. The three-cycle case has
lengths $u,v,v$, so its order is at most $uv\le n^2/8$.
This proves the quadratic bound used in the main theorem.

\section{The collision between two efficient pairs}\label{app:full-collision}

This appendix proves Lemma~\ref{lem:two-pairs}. Retain the parameters
$n=3m+2$, $p_0=(m+1)(2m+1)$, $v=2m(m+1)$, $b=p_0-1$, and $c=v-2$.

\begin{proof}
Normalize the first short-run entry to $(3,\ldots,n,2,1,0)$.
In a canonical $C$ cell, the nonzero cyclic core of its $r$th full run is
\[
 L_r=[\operatorname{rot}^{r-1}(1,3,\ldots,n-1,2),n],
 \qquad 1\le r\le n-3.
\]
Representatives beginning in $1$ give
\[
 L_1=[1,3,\ldots,n-1,2,n],\qquad
 L_{n-3}=[1,3,\ldots,n-3,n,n-2,n-1,2].
\]
Set $M=h_ch_b$ and $j=v-m-4$. Number cells from zero. The last full
run of cell $j$ in the second chain of the first pair has actual core
$h_b^{-1}g^{-j}L_{n-3}$. The first full run of cell $m$ in the second
chain of the second pair has core $M^{-1}h_b^{-1}g^{-m}L_1$.
Both cells precede their chains' final $B$ cell, since
$0\le j\le c-2$ and $m\le c-2$.

Equality of these two cores is equivalent to $EL_{n-3}=L_1$, where
\[
 E=g^mh_bh_cg^{-j}=g^m\rho g^{-1}\rho g^{m+2}.
\]
The last expression uses $g^{p_0}=1$, $c=p_0-m-3$ and
$j=p_0-2m-5$. The two nonzero cycles of $g$ are
\[
 (1,5,8,\ldots,3m+2),\qquad
 (2,4,7,\ldots,3m+1,3,6,\ldots,3m).
\]
Applying the displayed expression for $E$ gives
\[
\begin{gathered}
 E(1)=3,\quad E(2)=1,\quad E(x)=x+1\ (3\le x\le n-3),\\
 E(n-2)=2,\quad E(n-1)=n,\quad E(n)=n-1.
\end{gathered}
\]
It sends $L_{n-3}$ to $[3,4,\ldots,n-1,2,n,1]$, a rotation of $L_1$.
Thus the two full runs visit the same two-loop, contradicting class
distinctness.
\end{proof}

\subsection{The initial case on nine symbols}

At $m=2$, the efficient tiles are $(12)$, $(13,11)$ and $(14,10)$.
Their transfers lie in a group of order $12$. To see this directly, put
$A=(2\ 8)$, $Q=(2\ 8\ 7)$, $U=(1\ 3)$ and $V=(5\ 6)$.
The three transfers are $AV,AU,QUV$. They belong to
\[
 \{\sigma U^eV^f:
 \sigma\in\mathfrak S_{\{2,7,8\}},\ e,f\in\{0,1\},\
 \operatorname{sgn}(\sigma)=(-1)^{e+f}\},
\]
a group of order $12$. Twelve consecutive efficient interior tiles
would give thirteen prefix maps in this group, two equal, and hence
a repeated start class. Positive losses are integers and the final
allowance is two. Lemma~\ref{lem:tiles} therefore gives
$\CP(15,24-1/12)=\CP(15,287/12)$.


\section{The scalar conversion}\label{app:scalar}

For completeness, we give the counting step that turns a chain
inequality into a hereditary bound on segment deficits. This is Pantone's scalar theorem~\cite{pantone-framework}; no
integrality of $p$ or $t$ is required.

Call a finite sequence of nonnegative integers \emph{legal} if its
zero intervals have length at most $a$, except that an entirely zero
sequence may have length $a+1$; every factor $(0,1,d)$ has $d\ge2$;
and across an internal zero interval of length $a$, the neighboring
positive entries sum to at least $a+2$. These are the local run rules
for genuine light segments. A bound on all legal sequences satisfying
$\CP(p,t)$ therefore applies to those segments.

\begin{lemma}\label{lem:scalar}
Let $a\ge4$, $0\le p\le t$, $t>0$, and $t(a-2)\ge4p+6$.
Set
\[
 \lambda=at-2,\qquad B=2t(a+1)+4p.
\]
Every legal sequence satisfying $\CP(p,t)$ on its zero-ended factors
has length $r$ and sum $D$ satisfying $\lambda D-2tr+B\ge0$.
\end{lemma}

\begin{proof}
First consider a zero-ended sequence with $s\ge1$ positive clusters.
Let cluster $j$ have length $m_j$ and sum $D_j$. The rule concerning
$(0,1,d)$ gives $D_j\ge m_j+1$, so $D_j-2\ge m_j-1\ge0$.
Moreover $D_j=2$ holds exactly for a regular cluster. Put
\[
 M=\sum_j(m_j-1),\quad V=\sum_j(D_j-2),\quad
 b=\#\{j:(m_j,D_j)\ne(1,2)\}.
\]
Write $u_0,\ldots,u_s$ for the zero-gap lengths, including the two end
gaps, and define
\[
 g_0=\#\{j:1\le j<s,\ u_j=a\},\quad
 H=\sum_{j=1}^{s-1}\max(0,a-1-u_j),\quad
 L=(a-u_0)+(a-u_s).
\]
Counting entries and deficit gives
\begin{equation}\label{eq:cluster-count}
 r=as+a+1+M-H+g_0-L,\qquad D=2s+V.
\end{equation}

Let $h_1$ count the gaps of length $a-2$ between originally adjacent
regular clusters, and let $h_2$ count the shorter such gaps. Cutting
at the $b$ other clusters and at these $h_2$ gaps leaves at most
$b+1+h_2$ paths of chains. The sum of their numbers of links is $h_1$.
Thus $\CP(p,t)$ gives
\[
 s-b\le p(b+1+h_2)+\frac t2h_1.
\]
Since $H\ge h_1+2h_2$ and $p\le t$, it follows that
\begin{equation}\label{eq:J}
 J:=4p+4(p+1)b+2tH-4s\ge0.
\end{equation}

Set $c=\lambda-2t$. The numerical hypothesis says $c\ge4(p+1)$.
Give cluster $j$ the reserve
\[
 \Phi_j=\lambda(D_j-2)-2t(m_j-1)
       -4(p+1)\ind_{(m_j,D_j)\ne(1,2)}.
\]
A regular cluster has reserve zero. Every other cluster has
$D_j-2\ge1$, and therefore
$\Phi_j\ge c(D_j-2)-4(p+1)\ge0$.

Suppose a zero gap of length $a$ separates clusters $i,j$.
The actual neighboring positive entries sum to at least $a+2$;
each other entry in those clusters contributes at least one. Hence
\[
 (D_i-2)+(D_j-2)\ge a-2+(m_i-1)+(m_j-1).
\]
It follows that
\[
 \Phi_i+\Phi_j\ge\lambda(a-2)-8(p+1)
 \ge\lambda(a-2)-2c
 =\lambda(a-4)+4t\ge4t.
\]

Form a graph on the positive clusters, with an edge across each gap
of length $a$. Its $g_0$ edges form a subgraph of a path. In particular,
every degree is at most two. Substituting \eqref{eq:cluster-count}
and \eqref{eq:J} gives the identity
\[
 \lambda D-2tr+B
 =J+2tL
 +\frac12\sum_{\{i,j\}\in E}(\Phi_i+\Phi_j-4t)
 +\frac12\sum_j(2-\deg(j))\Phi_j.
\]
Every term on the right is nonnegative.

An entirely zero sequence has length at most $a+1$, so its left side
is at least $4p$. An entirely positive sequence satisfies
$\lambda D-2tr\ge0$. In a mixed sequence, remove the initial and
final positive clusters. Each removed entry contributes at least
$\lambda-2t=c\ge0$. The remainder is either one zero interval or a
zero-ended sequence of the kind already considered. The empty sequence
also satisfies the claimed inequality.
\end{proof}

\section{The hereditary forest estimate}\label{app:forest}

This appendix explains why the local deficit inequality applies to an
arbitrary covering word. A shortest word gives a minimum overlap path,
but that path need not visit every rotation class consecutively. The
forest construction replaces it by auxiliary complete-block paths and
records the cost of the replacement. The classes in different components
partition all the rotation classes. A block borrowed at the entrance to
one component is required to be new only to that component.

The construction uses a particular minimum path. Call a minimum path
\emph{reduced} if it has no edge from $u$ to its second cyclic rotation
$\sigma^2(u)$. Such a path exists. Among the minimum paths, choose one
with as many weight-one edges as possible. If it contained
$u\to\sigma^2(u)$, remove $\sigma(u)$ from its other location and
insert it between these two vertices. Removal cannot increase the weight,
by the triangle inequality, and insertion replaces one weight-two edge
by two weight-one edges. Moreover, removal deletes no weight-one edge.
The only possible weight-one predecessor of $\sigma(u)$ is $u$, whose
successor is already $\sigma^2(u)$, and the only possible weight-one
successor is $\sigma^2(u)$, whose predecessor is already $u$.
The new path would therefore still be minimum and have more weight-one
edges, a contradiction.

We state the properties of the forest construction for this reduced
minimum path. The construction and its residual inequality are inherited
from the lower-bound framework~\cite{pantone-framework}. The summation
argument after the statement is included here.

\begin{lemma}[Forest decomposition]\label{lem:forest-data}
Put $K=n+1\ge8$ and $T=n!$. A reduced minimum overlap path determines a finite
set of complete-block path components with a distinguished root. Each
nonroot component is classified as cheap or expensive. For component
$v$, let $b_v$ be its number of blocks, $q_v$ its number of runs,
$D_v$ its total deficit, $m_v$ its number of light segments, and
$d_v^-,d_v^+$ its first and last deficits. There are nonnegative
integers $A_v$ and $y_v$ such that
\[
 \sum_v b_v=T,\qquad b_v=nq_v-D_v,\qquad
 z_v=q_v+A_v,\qquad A_v\ge m_v-1.
\]
For the root and cheap components, $y_v=z_v$; for expensive components,
$y_v\ge z_v+1$.

For a cheap component with at least two runs, one can prepend an
entrance block to obtain another class-distinct path. Its block count
increases by one, its number of runs and $z_v$ are unchanged, and its
first deficit decreases by one. Its last deficit is unchanged.
A cheap component with one run has $y_v=1$ and
$d_v^-=d_v^+\ge1$.

There is a directed absorption relation on the components. Each component
has at most one incoming and one outgoing edge, and every target is a
cheap nonroot component. Every edge is either an isolated loop at a cheap
one-run component of deficit one, or an ordinary edge $v\to w$ with
\[
 d_v^+\ge2,\qquad d_w^-\ge2,\qquad d_v^++d_w^-\ge n+1.
\]
If $V$ is the number of cheap nonroot components with no incoming edge,
and $R=\sum_v y_v$, then
\begin{equation}\label{eq:residual}
 S(K)\ge K+(K+1)T-3+R+V.
\end{equation}
\end{lemma}

The last inequality accounts for the original path as well as the
auxiliary components. In particular, it does not assume that entrance
blocks borrowed by different components are distinct. This is the reason
for proving the local inequality on every genuine segment, including
consecutive factors of an entrance-extended path.

\begin{theorem}[Hereditary forest estimate]\label{thm:forest}
Suppose $0<\delta<1$ and $1\le\beta\le\alpha:=n-\delta$.
If every genuine light segment with $r$ runs and total deficit $D$
satisfies $D\ge\delta r-\beta$, then
\[
 S(K)\ge K+(K+1)n!-3+
 \ceil{\frac{n!-\beta}{n-\delta}}.
\]
\end{theorem}

\begin{proof}
The first step recovers endpoint information from the hypothesis.
Write $F(d)=\max(0,d-\delta)$. Deleting the first deficit from a
segment and applying the hypothesis to the remainder gives an extra
term $d^- -\delta$. Applying the hypothesis without deletion gives
an extra term zero. Taking the larger of these two estimates gives
\[
 D\ge\delta r-\beta+F(d^-).
\]
There is an identical estimate at the last endpoint. When $r\ge2$,
we can delete either endpoint, both, or neither, and obtain
\begin{equation}\label{eq:two-endpoints}
 D\ge\delta r-\beta+F(d^-)+F(d^+).
\end{equation}
The empty remainder causes no difficulty because $0\ge-\beta$.
The two endpoints in \eqref{eq:two-endpoints} must be distinct;
we shall handle one-run components separately.

Apply Lemma~\ref{lem:forest-data} and put
$s_v=b_v-\alpha y_v$. Summing the segment estimates within one
component gives
\[
 b_v-\alpha z_v
 \le m_v\beta-\alpha A_v-F(d_v^+)
 \le\beta-F(d_v^+),
\]
because $A_v\ge m_v-1$ and $\alpha\ge\beta$. Thus
\[
 s_v\le\beta-F(d_v^+)\quad\hbox{at the root},\qquad
 s_v\le\beta-\alpha-F(d_v^+)\quad\hbox{at an expensive component}.
\]
For a cheap component with at least two runs, apply the same argument
to its entrance-extended path and use both endpoint estimates. Define
$A(d)=\max(0,d-1-\delta)$ and $c=\beta-1$. The additional block gives
\[
 s_v\le c-A(d_v^-)-F(d_v^+).
\]
If the two ends lie in different light segments, apply the one-end
estimate to each of those segments. If they lie in the same segment,
that segment has at least two runs and \eqref{eq:two-endpoints} applies.
For a cheap component with a single run, the exact identity is
\[
 s_v=(n-d_v^-)-(n-\delta)=\delta-d_v^-=-F(d_v^-).
\]

Put $\widehat A(d)=\min(A(d),c)$. For a cheap component, let
$i_v,o_v\in\{0,1\}$ indicate whether it has an incoming and an
outgoing absorption edge. Then both the multi-run and single-run
estimates imply
\[
 s_v\le c-i_v\widehat A(d_v^-)-o_vF(d_v^+).
\]
For a single run, this follows from $0\le\widehat A\le c$.
Clipping the incoming charge is what permits one formula to cover that
case without counting its only deficit twice.

Assign an intercept $c_v$ to each component: use $\beta$ at the root,
$c$ at cheap nonroots, and $\beta-\alpha$ at expensive nonroots.
Assign the incoming charge $a_v=\widehat A(d_v^-)$ at cheap vertices
and zero elsewhere, and the outgoing charge $b'_v=F(d_v^+)$.
There is one exception. At a vertex carrying an isolated special loop,
use $c_v=\delta-1$ and $a_v=b'_v=0$; its surplus is exactly
$\delta-1$.

The preceding bounds give
\[
 s_v+i_va_v+o_vb'_v\le c_v.
\]
For an ordinary edge $v\to w$, its two deficits give
\[
 A(d_w^-)+F(d_v^+)\ge n-2\delta,\qquad
 F(d_v^+)\ge2-\delta.
\]
Consequently its charges cover the target intercept, since
\[
 b'_v+a_w-c_w
 =\min\{F(d_v^+)+A(d_w^-)-c,\ F(d_v^+)\}
 \ge\min\{n+1-\beta-2\delta,\ 2-\delta\}>0.
\]
For a special loop, the corresponding excess is $1-\delta>0$.

Now sum the vertex inequalities and the edge inequalities. Each incoming
and outgoing charge occurs exactly once on either side and cancels.
The remaining intercepts belong to vertices with no incoming edge.
These are the root, the $V$ unabsorbed cheap vertices, and all the
expensive vertices. The latter have nonpositive intercept. It follows that
\[
 T-\alpha R=\sum_v s_v\le\beta+(\beta-1)V.
\]
Since $\beta\le\alpha$, this implies
$T\le\alpha(R+V)+\beta$. The number $R+V$ is an integer, so
\eqref{eq:residual} gives the asserted ceiling bound.
\end{proof}

\subsection{From the chain estimate to word length}

\begin{proof}[Proof of Proposition~\ref{prop:transfer}]
Set
\[
 \lambda=(n-2)t-2,\qquad
 \delta=\frac{2t}{\lambda},\qquad
 \beta=\frac{2t(n-1)+4p}{\lambda}.
\]
Lemma~\ref{lem:scalar} gives
$D\ge\delta r-\beta$ for the total deficit $D$ of every segment with
$r$ runs. The same estimate holds on its consecutive run factors.
The parameter assumptions give $0<\delta<1$ and $1<\beta<n-\delta$.
For the latter inequality, use $4p\le t(n-4)-6$ and
$t(n-4)\ge6$ to obtain
\[
 n(n-4)t-4p-2n\ge (n-1)(n-4)t+6-2n\ge4n>0.
\]
Theorem~\ref{thm:forest}
therefore gives
\begin{equation}\label{eq:forest}
 S(K)\ge K+(K+1)n!-3+
 \ceil{\frac{n!-\beta}{n-\delta}}.
\end{equation}
Substitution and removal of the integer $(n-1)!$ from the ceiling give
\eqref{eq:bound}.
\end{proof}

\section{Parameter checks and integer improvements}\label{app:numerical}

\subsection{The transfer hypotheses}

The chain inequalities are \eqref{eq:CP1}, \eqref{eq:CP2}, and
Proposition~\ref{prop:CP0}. We verify the numerical assumptions of
Proposition~\ref{prop:transfer}.

For $n=3m\ge9$, write $p_0=2m$. Since $R_0\ge1$, the new parameters
satisfy
\[
 t-p=p_0-\frac72-\frac3{4R_0}>0,
\]
and
\[
 t(n-4)-4p-6
 \ge(2p_0-3)(n-4)-4p_0-8-\frac{n-2}{2}
 =\frac43n^2-\frac{23}{2}n+5>0.
\]
The last polynomial is positive and increasing for $n\ge9$.

For $n=3m+1$ with $m\ge3$, the new $t$ is at least $t_0-1\ge p_0$,
and
\[
 (t_0-1)(n-4)-4p_0-6=6m^3-8m^2-19m+3>0.
\]
Writing $m=3+x$ makes the right side
$6x^3+46x^2+95x+36$.
For $n=3m+2$, $m\ge2$, all the stated choices have $t\ge t_0-1$.
Here
\[
 t_0-1-p_0=2(m+2)(m-1)>0,
\]
and
\[
 (t_0-1)(n-4)-4p_0-6
 =(m-2)(12m^2+23m+15)+26>0.
\]
At $K=8$, the earlier chain theorem gives $\CP(10,10)$
\cite{pantone-framework}; increasing $t$ to $46/3$ preserves this
inequality and gives $t(n-4)=4p+6$. The values for $K=9,12$ are the
two initial cases in Section~\ref{sec:pairs}. Those for $K=10,11$
are the first instances of the other two residue formulas.

Proposition~\ref{prop:transfer} now proves the asserted word bound.
Replacing either return order by any larger real number preserves the
counts in Lemma~\ref{lem:tiles} or \eqref{eq:alternating-count}.
The weaker quadratic choices still satisfy the same numerical
estimates, so they too give the claimed bounds.

\subsection{Exact changes in the bound}

Put $T=n!$ and $\kappa=n(n-2)-2$. The nonconstant real term in
\eqref{eq:forest} is
\[
 f(p,t)=\frac{t((n-2)T-2(n-1))-2T-4p}{\kappa t-2n}.
\]
First suppose $n=3m+1$. With $h=1/R_1$, direct subtraction gives
\begin{equation}\label{eq:gain1}
 f(p_0,t_0-h)-f(p_0,t_0)
 =\frac{4h[T-p_0\kappa-n(n-1)]}
 {(\kappa(t_0-h)-2n)(\kappa t_0-2n)}.
\end{equation}
For $n\ge10$, we have $t_0<n^2/4$, $p_0<n^2$,
$\kappa<n^2$, and $h\ge4/n^2$. Hence this difference exceeds
\[
 256\frac{n!-n^4-n^2}{n^{10}}.
\]
At $n=13$ the expression exceeds one, and it increases thereafter:
$n!/n^{10}$ increases since $(1+1/n)^9<3<n$, while
$n^{-6}$ and $n^{-8}$ decrease. Thus the integer ceiling improves
at every stated size in this residue class. Since
$t_0\sim2n^2/9$ and $\kappa\sim n^2$, \eqref{eq:gain1} is
asymptotic to $81n!/(R_1n^8)$. The quadratic upper bound on $R_1$
gives the constant $324$.

For $n=3m$, set $h=1/(2R_0)$, $d=1/2+h/2$, and
$D_0=\kappa t_0-2n$. Subtraction now gives
\[
 f(p_0+d,t_0-h)-f(p_0,t_0)
 =\frac{4h[T-n(n-1)-\kappa p_0-(R_0+1/2)D_0]}
 {D_0(D_0-\kappa h)}.
\]
Using $R_0\le n^2/4$, $p_0=2n/3$, $t_0<4n/3$ and
$\kappa<n^2$, the total subtracted from $T$ is less than $n^5$.
The difference therefore exceeds
\[
 \frac92\frac{n!-n^5}{n^8}.
\]
This is greater than one at $n=12$ and increases thereafter.
Its factorial term increases because $(1+1/n)^7<3<n$; the
subtracted term decreases. More precisely, the exact difference
is asymptotic to $9n!/(8R_0n^6)$, giving the constant $9/2$.
The difference of two ceilings changes an unrounded difference by
less than one. Finally, $n!/n^j\sim(n-j)!$, and $K=n+1$,
which gives both normalizations in the statement.

\section{The formal statements}

The companion Lean development separates its public definitions from
the proofs. The coverage predicate on a word $w$ over $\operatorname{Fin}K$
is the assertion that every list $p$ of length $K$ with distinct entries
has a prefix $u$ and a suffix $v$ with $w=upv$. Thus the final statements
refer directly to covering words, with no path or chain hypothesis.

The two general statements in
\nolinkurl{Submissions/Challenge.lean}
are Theorem~\ref{thm:main} with the exact and quadratic return orders.
Their formulas include the initial rows $8\le K\le12$.
Their declaration names are \nolinkurl{exact_word_bound} and
\nolinkurl{quadratic_word_bound}, in the namespace
\nolinkurl{SuperpermutationPalomar.GeneralLower}.
The challenge imports only Mathlib. The solution applies the original
theorems to these statements, with Lean checking that their definitions
agree.
The Lean formula uses the natural-number ceiling. The correction in
\eqref{eq:bound} is positive in every stated application, so this
agrees with the ordinary ceiling used here.

The same challenge states thirteen asymptotic consequences. The correspondence
with the paper is recorded below. Names in the second column are the
suffixes of declarations in
\nolinkurl{SuperpermutationPalomar.LowerAsymptotics}.

\begin{center}
\small
\begin{tabular}{@{}>{\raggedright\arraybackslash}p{0.37\textwidth}
                  >{\raggedright\arraybackslash}p{0.59\textwidth}@{}}
\toprule
Mathematical statement & Lean declaration\\
\midrule
The leading correction in Corollary~\ref{cor:comparison}
& \nolinkurl{exact_leading_limit}\newline
  \nolinkurl{quadratic_leading_limit}\\[3pt]
Corollary~\ref{cor:word-asymptotic}
& \nolinkurl{word_lower_bound_two_sub_epsilon}\\[3pt]
The three comparisons of $B$ with $H$
& \nolinkurl{exact_reference_limit_zero}\newline
  \nolinkurl{exact_reference_limit_one}\newline
  \nolinkurl{exact_reference_limit_two}\\[3pt]
The three comparisons of $Q$ with $H$
& \nolinkurl{quadratic_reference_limit_zero}\newline
  \nolinkurl{quadratic_reference_limit_one}\newline
  \nolinkurl{quadratic_reference_limit_two}\\[3pt]
The two limits for $Q-O$
& \nolinkurl{quadratic_gain_zero_limit}\newline
  \nolinkurl{quadratic_gain_one_limit}\\[3pt]
The two eventual lower estimates for $B-O$
& \nolinkurl{exact_gain_zero_asymptotic_lower}\newline
  \nolinkurl{exact_gain_one_asymptotic_lower}\\
\bottomrule
\end{tabular}
\end{center}

Here the suffixes zero, one and two refer to $n=K-1$ modulo three.
The last two declarations say that, for each positive epsilon, the
normalized difference is eventually at least the stated constant minus
that epsilon. They do not assert convergence. The finite strict-improvement
thresholds in Proposition~\ref{prop:gains} have the ordinary proof in
Appendix~\ref{app:numerical}; they are not separate claims in this
submission.

Every declaration just listed, and both general word bounds, has a proof
using only \nolinkurl{propext}, \nolinkurl{Classical.choice}, and
\nolinkurl{Quot.sound}. No native computation assumption is used in these
lower-bound families.

\begin{thebibliography}{9}
\raggedright
\bibitem{pantone-framework}
J. Pantone, \emph{The lower-bound framework for superpermutations},
companion Lean library, namespace \nolinkurl{SuperpermutationLowerBound}.
The paths below are relative to
\nolinkurl{proofs/SuperpermutationLowerBound/}.
The geometric inputs are in \nolinkurl{Refinement/ReturnGeometry.lean},
\nolinkurl{Refinement/ReturnLocalCollisions.lean}, and the
\nolinkurl{Refinement/ChainPacking.lean} family of modules.
The forest reduction is in the \nolinkurl{Section01}--\nolinkurl{Section10}
directories and \nolinkurl{Refinement/ForestResidual.lean}; the hereditary
assembly is in \nolinkurl{Refinement/ForestAssembly.lean}.

\bibitem{halfar}
J. Halfar, \emph{Three lower bounds in Lean: Theorems A, B and C},
\href{https://github.com/jhalfar/superpermutations/tree/cc5c763bd5854317b20dadb138990802656e6f03/proofs/lean/lower}{lower-bound proof source},
revision \nolinkurl{cc5c763bd5854317b20dadb138990802656e6f03},
October 2026.
\bibitem{halfar-announcement}
J. Halfar, \emph{New upper bound with 1771/3456 and new lower bounds
for $n=8$ to $14$, both in Lean}, announcement to the Superpermutators
group, 6 October 2026.
\url{https://groups.google.com/g/superpermutators/c/6ozhk3faHDA}.
\end{thebibliography}

\end{document}
